% =====================================================================
%  Thème 5 — Résolution des triangles — NIVEAU MOYEN
%  Exercices
% =====================================================================
\documentclass[11pt,a4paper]{article}
\usepackage[utf8]{inputenc}
\usepackage[T1]{fontenc}
\usepackage{lmodern}
\usepackage[french]{babel}
\usepackage{amsmath,amssymb,mathtools}
\usepackage{icomma}
\usepackage{xcolor}
\usepackage{colortbl}
\usepackage{tikz}
\usepackage[most]{tcolorbox}
\usepackage{enumitem}
\usepackage{array}
\usepackage[a4paper,margin=2.2cm]{geometry}
\usetikzlibrary{arrows.meta,calc}

\definecolor{primary}{RGB}{214,51,108}
\definecolor{primarydark}{RGB}{140,20,64}

\DeclareMathOperator{\tg}{tg}
\DeclareMathOperator{\cotg}{cotg}
\newcommand{\degr}{^{\circ}}
\newcommand{\Z}{\mathbb{Z}}

\newtcolorbox{consigne}{enhanced,breakable,colback=primary!5,colframe=primary,
  boxrule=0.8pt,arc=2pt,left=3mm,right=3mm,top=2mm,bottom=2mm}

\newcounter{exo}
\newcommand{\exo}[1]{\medskip\refstepcounter{exo}%
  \noindent{\large\bfseries\color{primarydark}Exercice \theexo}%
  \ \textcolor{primary}{\textbullet}\ \textbf{#1}\par\nopagebreak\smallskip}

\setlength{\parindent}{0pt}
\pagestyle{empty}

\begin{document}

\begin{center}
{\LARGE\bfseries\color{primary}Thème 5 — Résolution des triangles}\\[2pt]
{\color{primarydark}\textsc{Niveau Moyen — Exercices}}
\end{center}
\hrule height 1pt
\medskip

\begin{consigne}
\textbf{Objectif :} enchaîner \emph{deux} étapes (aire, loi des sinus, loi des cosinus, SOH-CAH-TOA)
et identifier un angle à partir d'une relation entre côtés. \textbf{Sans calculatrice}, valeurs
\emph{exactes}. Convention : $a=\overline{BC}$, $b=\overline{CA}$, $c=\overline{AB}$.
\end{consigne}

\exo{De l'aire au troisième côté}
Dans le triangle $ABC$, on sait que $\hat C=60\degr$, $a=\overline{BC}=8$ et l'aire vaut
$\mathcal{A}=8\sqrt3$.
\begin{enumerate}[label=\textbf{(\alph*)},leftmargin=2em,itemsep=2pt]
\item À l'aide de la formule de l'aire, détermine le côté $b=\overline{CA}$.
\item En déduis $c=\overline{AB}$ par la loi des cosinus.
\end{enumerate}

\exo{Loi des sinus puis loi des cosinus}
Un triangle $ABC$ vérifie $\hat A=\hat B=30\degr$ et $a=\overline{BC}=5$.
\begin{enumerate}[label=\textbf{(\alph*)},leftmargin=2em,itemsep=2pt]
\item Calcule l'angle $\hat C$.
\item Détermine le côté $b=\overline{CA}$ (loi des sinus). Que peux-tu dire du triangle ?
\item Calcule $c=\overline{AB}$ par la loi des cosinus, puis vérifie le résultat avec la loi des sinus.
\end{enumerate}

\exo{Un triangle rectangle à résoudre entièrement}
Le triangle $ABC$ est rectangle en $A$ ; on donne $\hat B=60\degr$ et $c=\overline{AB}=4$.
\begin{enumerate}[label=\textbf{(\alph*)},leftmargin=2em,itemsep=2pt]
\item Calcule le côté $\overline{AC}$ (SOH-CAH-TOA) puis l'hypoténuse $\overline{BC}$.
\item Vérifie ta réponse avec le théorème de Pythagore, puis calcule l'aire du triangle.
\end{enumerate}

\exo{Reconnaître un angle à partir d'une relation entre côtés}
\begin{enumerate}[label=\textbf{(\alph*)},leftmargin=2em,itemsep=2pt]
\item Un triangle vérifie $c^2=a^2+b^2-ab$. Que vaut l'angle $\hat C$ ?
\item Un triangle vérifie $c^2=a^2+b^2+ab$. Que vaut l'angle $\hat C$ ?
\item Les côtés d'un triangle mesurent $7$, $8$ et $13$. Calcule son plus grand angle
      (il est opposé au plus grand côté).
\end{enumerate}

\exo{Aire et diagonale d'un parallélogramme}
$ABCD$ est un parallélogramme avec $\overline{AB}=6$, $\overline{AD}=4$ et l'angle $\hat A=60\degr$.
\begin{enumerate}[label=\textbf{(\alph*)},leftmargin=2em,itemsep=2pt]
\item Calcule l'aire du parallélogramme (rappel : aire $=$ produit de deux côtés adjacents
      $\times\sin$ de l'angle entre eux).
\item Calcule la longueur de la diagonale $\overline{BD}$ (loi des cosinus dans le triangle $ABD$).
\end{enumerate}

\end{document}
